\chapter{RC610 line printer}\label{rc610-line-printer}
\index{RC610}\index{Printer}

All supplied monitor profiles include the external process \texttt{printer},
on device 7, interrupt 10. This applies to compatibility, protected, MIPS/TS
and DFC403 operation. There is no printer-free maintained monitor variant.
The native emulator writes to an ordinary host file; it does not submit jobs
to a host printing service.

\section{Automatic output and file selection}

The supplied configurations open output automatically in append mode:

\begin{center}
\begin{tabular}{ll}
\textbf{Configuration} & \textbf{Output file} \\
\texttt{RC4000.cfg} & \texttt{printer.txt} \\
\texttt{RC4000.protected.cfg} & \texttt{printer.protected.txt} \\
\texttt{RC4000.mipsts.cfg} & \texttt{printer.mipsts.txt} \\
\texttt{RC4000.dfc403.cfg} & \texttt{printer.dfc403.txt} \\
\end{tabular}
\end{center}

Paths are relative to the configuration file. Existing contents survive
restarts; the emulator appends new prints and never truncates an output file.
Each profile has its own output file. Use separate working directories for
concurrent sessions; appending does not provide guest job scheduling or keep
two emulators' jobs separate.
The DFC403 configuration is an additional source/build-tree profile; the
binary archive supplies the other three configurations.

\begin{commandexample}{Configuration: automatic append-mode printer}
DEFINE RC610 DEVNO=7 INTNO=10 FILE=printer.txt APPEND=YES
\end{commandexample}

\texttt{APPEND=NO}, or omitting \texttt{APPEND}, requires a nonexistent output
file. Omitting \texttt{FILE} leaves the device installed but offline until an
output file is attached. Only regular files are accepted; symbolic links,
FIFOs and device files are rejected. New files have private host permissions.

To change the destination, suspend execution with Control-E and enter:

\begin{commandexample}{Emulator CLI: choose a new exclusive output file}
show media
detach 7
attach 7 session-print.txt
continue
\end{commandexample}

For an existing file use \texttt{attach 7 session-print.txt append}. Bare
\texttt{attach} never overwrites an existing file. These CLI paths are relative
to the emulator's working directory, unlike configuration paths.
\texttt{show media} reports printer status, counts and its output path.

Wait for guest output to finish before detaching. An active hardware print
cannot be detached: resume execution, let it complete, then try again.
An incomplete line without a printing control is discarded and counted on
close, not retroactively printed. CPU reset preserves the device and file.
Normal shutdown completes an already accepted hardware print and closes output;
abnormal termination is not crash-safe spooling.

\section{A first print from s}

After date initialization, request RC attention and select \texttt{s}.
The following bounded memory dump prints only the monitor's IRQ-0 vector:

\begin{commandexample}{Guest system process s: print a known monitor word}
print 16 16
\end{commandexample}

Wait for \texttt{ready}. The host printer file gains one formatted line.
Repeating the command prints another line; the values depend on the monitor
build. This has been checked with compatibility and protected monitors and
with the DFC403 monitor. It is a memory-dump command, not file printing.
Do not substitute an arbitrary user-memory range and expect a source listing.

\section{Print an installed Algol source}

The supplied disk contains \texttt{cpuregtest}, the Algol source used by
\texttt{regnetest}. The tutorial example \texttt{examples/print\_source.alg}
opens that guest file for input and the external printer for output, copies
characters up to the guest text end marker, and closes both zones. It prints
the source; it does not compile or execute the numerical test itself.

Copy the example into the configured \texttt{exchange} directory as
\texttt{print\_source.alg}. In \texttt{s}, allocate a temporary FP:

\begin{commandexample}{Guest system process s: create a printing session}
new ptest size 100000 temp disc 200 10 run
\end{commandexample}

\begin{commandexample}{Guest FP, compatibility profile: compile and print}
psrc=fromhost "print_source.alg"
ptext=algol psrc connect.no
ptext
\end{commandexample}

For the protected profile, start with \texttt{RC4000.protected.cfg} and its
matching monitor/disk pair. Use the same example and \texttt{s} allocation,
but import through the protected profile's transfer program:

\begin{commandexample}{Guest FP, protected profile: compile and print}
psrc=rcin "print_source.alg"
ptext=algol psrc connect.no
ptext
\end{commandexample}

Both checked profiles print 630 lines/printing operations and 17,309
bytes, beginning with the MPROFILE header. The output is visible before closing
the emulator. The tiny program uses \texttt{open(source,4,...)} for the disk
area and \texttt{open(printer,0,...)} for the physical printer; selecting a
disk output called \texttt{printer} is not the same operation.
The FP recipes \texttt{o printer} and \texttt{printer=copy cpuregtest} are
not accepted substitutes: exploratory runs did not drive the hardware printer.

Both runs match an independent whole-source guest export: \texttt{tohost} in
compatibility mode, \texttt{rcout} in protected mode. The protected check keeps
\texttt{PROTECTION MODE=RC4000} enabled and observes keyed task instructions
during execution of the printing program, not merely during boot. It uses
the existing protected libraries without changes to the example or images.
Enter \texttt{finis} after printing and wait for \texttt{pause ptest finis fp}.
Normal return to \texttt{s} and a subsequent \texttt{list} are also checked.

\section{Text mapping, visibility and limits}

The host file is a print-operation stream, not a rendered paper image. A
printing control emits the buffered line followed by its original control:

\begin{center}
\begin{tabular}{lll}
\textbf{Guest control} & \textbf{Host byte} & \textbf{Meaning} \\
NL (10) & LF & Begin the next line. \\
VT (11) & VT & Retain vertical-tab motion. \\
FF (12) & FF & Retain the form/page separator. \\
CR (13) & CR & Return on the same line for overprinting. \\
\end{tabular}
\end{center}

No LF is inserted after CR. VT/FF remain literal controls; no physical
format tape is emulated. Text viewers may therefore display overprints and
page controls differently. The buffer has 132 columns; excess graphics are
lost and counted. Unsupported standard-character-set codes and optional
HT/SO/SI/CAN/ESC functions are ignored. National RC character positions remain
raw seven-bit bytes, not UTF-8 Danish letters.

Every completed print calls \texttt{fflush} before its completion interrupt,
including CR, VT and FF. Output does not wait for a full host buffer or emulator
exit. A guest program can still buffer text before sending it to the printer.
There is no per-line \texttt{fsync} durability guarantee. Host write/flush errors
are latched as virtual offline/paper-out status, not silently ignored or
replayed; choose a healthy new output after resolving the host error.

The timing model is deterministic, not mechanically accurate. A print completes
on the next positive device advance. The supplied MIPS/TS profile supports
PRIMO transports to the real external process \texttt{printer}; see
\manualref{mips-ts-primo-printing}{Print through PRIMO}. Repeated jobs,
signoff and restart append are checked. Recovery of a transport submitted while
output is detached is described in
\manualref{mips-ts-printer-recovery}{Recover a job submitted with output detached}.
General host-write failure recovery and arbitrary queue capacity remain separate
tasks. The MIPS/TS chapter records the bounded two-client check and its second
blocking-wait restriction.
